\chapter{Numerical and Statistical Stability}

\section{Ill-Conditioning}

In previous chapters, we studied optimization methods and their practical behavior. However, the effectiveness of these methods depends critically on the numerical properties of the problem.

\medskip

\noindent
\textbf{Definition (informal).}  
A problem is said to be \emph{ill-conditioned} if small changes in the input lead to large changes in the output.

\medskip

\noindent
\textbf{Example (linear regression).}  
The solution
\[
w^\star = (X^\top X)^{-1} X^\top y
\]
is sensitive to perturbations when $X^\top X$ is nearly singular.

\medskip

\noindent
\textbf{Condition number.}  
For a positive definite matrix $A$, the condition number is
\[
\kappa(A) = \frac{\lambda_{\max}(A)}{\lambda_{\min}(A)}.
\]

\medskip

\noindent
\textbf{Interpretation.}
\begin{itemize}
    \item Large $\kappa(A)$ $\rightarrow$ ill-conditioned problem
    \item Small $\kappa(A)$ $\rightarrow$ well-conditioned problem
\end{itemize}

\medskip

\noindent
\textbf{Geometric intuition.}  
Ill-conditioning corresponds to elongated level sets of the loss function, making optimization difficult.

\medskip

\noindent
\textbf{Consequences.}
\begin{itemize}
    \item Slow convergence of gradient methods
    \item Numerical instability in matrix inversion
    \item Sensitivity to noise
\end{itemize}

\medskip

\noindent
\textbf{Sources of ill-conditioning.}
\begin{itemize}
    \item Highly correlated features
    \item Redundant inputs
    \item Poor scaling of variables
\end{itemize}

\section{Scaling and Normalization}

A simple but effective way to improve conditioning is to preprocess the data.

\medskip

\noindent
\textbf{Feature scaling.}  
Transform each feature to have comparable magnitude.

\medskip

\noindent
\textbf{Standardization.}
\[
x_j \leftarrow \frac{x_j - \mu_j}{\sigma_j},
\]
where $\mu_j$ and $\sigma_j$ are the mean and standard deviation of feature $j$.

\medskip

\noindent
\textbf{Normalization.}  
Scale inputs to have bounded norms, e.g.,
\[
x \leftarrow \frac{x}{\|x\|}.
\]

\medskip

\noindent
\textbf{Effect on optimization.}
\begin{itemize}
    \item Makes gradients more balanced across dimensions
    \item Improves convergence speed
    \item Reduces numerical instability
\end{itemize}

\medskip

\noindent
\textbf{Connection to conditioning.}  
Proper scaling reduces the condition number of the problem.

\medskip

\noindent
\textbf{Normalization in deep learning.}
\begin{itemize}
    \item Batch normalization
    \item Layer normalization
\end{itemize}

\medskip

\noindent
\textbf{Insight.}  
Normalization can be viewed as dynamically adjusting the geometry of the optimization problem.

\section{Sensitivity to Noise}

Real-world data is often noisy. A stable learning algorithm should not be overly sensitive to small perturbations.

\medskip

\noindent
\textbf{Noise sources.}
\begin{itemize}
    \item Measurement errors
    \item Incomplete observations
    \item Intrinsic randomness
\end{itemize}

\medskip

\noindent
\textbf{Effect on learning.}
\begin{itemize}
    \item Can lead to overfitting
    \item Can destabilize parameter estimates
\end{itemize}

\medskip

\noindent
\textbf{Example (least squares).}  
In ill-conditioned problems, small noise in $y$ can cause large changes in $w^\star$.

\medskip

\noindent
\textbf{Bias--variance perspective.}
\begin{itemize}
    \item High-variance models are more sensitive to noise
    \item Regularization reduces sensitivity
\end{itemize}

\medskip

\noindent
\textbf{Stability as robustness.}  
A stable algorithm produces similar outputs when inputs are slightly perturbed.

\medskip

\noindent
\textbf{Algorithmic stability (informal).}  
If removing or modifying a single data point does not significantly change the learned model, the algorithm is considered stable.

\medskip

\noindent
\textbf{Insight.}  
Stability is closely related to generalization: stable algorithms tend to generalize better.

\section{Robustness Considerations}

Beyond numerical stability, we are often concerned with robustness to broader forms of perturbation.

\medskip

\noindent
\textbf{Types of robustness.}
\begin{itemize}
    \item Robustness to noise in data
    \item Robustness to distribution shifts
    \item Robustness to adversarial perturbations
\end{itemize}

\medskip

\noindent
\textbf{Robust loss functions.}
\begin{itemize}
    \item Mean squared error: sensitive to outliers
    \item Absolute error: more robust
\end{itemize}

\medskip

\noindent
\textbf{Regularization and robustness.}  
Constraining model complexity often improves robustness.

\medskip

\noindent
\textbf{Data augmentation.}  
Introducing variations in training data can improve robustness to transformations.

\medskip

\noindent
\textbf{Adversarial perspective (informal).}  
Small, carefully chosen perturbations can significantly affect model predictions, especially in high-dimensional settings.

\medskip

\noindent
\textbf{Insight.}  
Robustness requires both stable models and appropriate assumptions about the data.

\section{A Unifying Perspective}

Numerical and statistical stability are closely related:

\begin{itemize}
    \item \textbf{Numerical stability:} sensitivity to computational errors
    \item \textbf{Statistical stability:} sensitivity to data variations
\end{itemize}

\medskip

\noindent
\textbf{Key factors.}
\begin{itemize}
    \item Conditioning of the problem
    \item Data preprocessing
    \item Model complexity
    \item Regularization
\end{itemize}

\medskip

\noindent
\textbf{Core insight.}  
Stable learning systems are those that behave predictably under small perturbations, both numerical and statistical.

\section{Outlook}

This chapter examined stability from both numerical and statistical perspectives:
\begin{itemize}
    \item ill-conditioning and its consequences,
    \item the role of scaling and normalization,
    \item sensitivity to noise,
    \item robustness considerations.
\end{itemize}

\medskip

In the next chapter, we transition from classical models to neural networks, motivating deep learning as a natural extension of these ideas.

\medskip

\noindent
\textbf{Guiding principle.}  
Stability is essential for reliable learning: models should not only fit data, but do so in a way that is robust to perturbations and uncertainty.